% Source PDF page 88; manual page 2-59.
\subsection{Propulsive Forces}\label{sec:propulsive-forces}

The magnitude and direction of each propulsive force vector are specified in the
table file, as described in \cref{sec:table-types}, or directly in the problem file,
as described in \cref{sec:block-prop}. If more than one force vector is supplied,
the vectors are summed to form $\sum\vec F_{\mathrm{prop}}$.

Because the equations of motion treat the vehicle as a point mass, every thrust
vector acts through the vehicle center of mass. Its direction is determined by thrust
vector angles $\epsilon_1$ and $\epsilon_2$, shown in
\cref{fig:chapter2-thrust-vector-angles}.

\begin{figure}[htbp]
  \centering
  \resizebox{0.86\linewidth}{!}{\begin{tikzpicture}[>=Latex,line cap=round,line join=round,scale=0.88]
  \coordinate (o) at (-3.6,0);
  \coordinate (p) at (2.2,1.45);
  \draw[->,line width=1.0pt] (o) -- (3.5,0) node[right] {$\hat x_b$};
  \draw[->,line width=1.0pt] (o) -- (-3.6,-1.35) node[below] {$\hat z_b$};
  \draw[->,line width=1.0pt] (o) -- (-2.5,-0.95) node[below right] {$\hat y_b$};
  \node[above left] at (o) {vehicle};
  \node[above left] at (-3.0,1.05) {$\hat y_b$ is out of the page};

  % Thrust vector and its projection.
  \draw[<-,line width=1.15pt] (o) -- (p)
    node[midway,above left] {$\vec F_{\mathrm{prop}}$};
  \draw[dashed] (p) -- (2.2,0);
  \draw[->,line width=0.85pt] (o) ++(1.45,0) arc[start angle=0,end angle=14,radius=1.45]
    node[midway,right] {$\epsilon_1$};

  % yz-plane at the projected point.
  \draw (2.2,0) ellipse (0.58 and 1.35);
  \draw[->] (2.2,0) -- (2.2,-1.55) node[below] {$\hat z_b$};
  \draw[->] (2.2,0) -- (3.15,-0.75) node[right] {$\hat y_b$};
  \draw[line width=0.9pt] (2.2,0) -- (2.55,0.92);
  \draw[->,line width=0.8pt] (2.2,0.62) arc[start angle=90,end angle=69,radius=0.62]
    node[midway,right] {$\epsilon_2$};
  \node[align=left,right] at (2.85,1.35)
    {projection of $\vec F_{\mathrm{prop}}$\\onto the $\hat y_b$--$\hat z_b$ plane};
\end{tikzpicture}
}
  \caption{Thrust vector angles.}
  \label{fig:chapter2-thrust-vector-angles}
\end{figure}

The angle $\epsilon_1$ is analogous to total angle of attack: it is the angle between
the body $x$ axis and the thrust vector and can range from $0^\circ$ to
$180^\circ$. The angle $\epsilon_2$ is analogous to windward meridian: it is the
angle between the projection of the thrust vector onto the $\hat y_b$--$\hat z_b$
plane and the negative $\hat y_b$ axis.

The body-frame thrust vector is
\begin{equation}
  \begin{aligned}
  \vec F_{\mathrm{prop}}
  ={}&\lVert\vec F_{\mathrm{prop}}\rVert\cos\epsilon_1\,\hat x_b\\
    &-\lVert\vec F_{\mathrm{prop}}\rVert
      \sin\epsilon_1\cos\epsilon_2\,\hat y_b\\
    &-\lVert\vec F_{\mathrm{prop}}\rVert
      \sin\epsilon_1\sin\epsilon_2\,\hat z_b.
  \end{aligned}
  \label{eq:propulsive-force-vector}
\end{equation}
The function \texttt{force\_prop} evaluates this expression and transforms the
result from body coordinates to ECFC coordinates using the method in
\cref{sec:coordinate-transformations}.

The associated output variables are
\begin{center}
\begin{tabularx}{0.78\linewidth}{@{}>{\ttfamily}l>{$}l<{$}X@{}}
  thrust & \lVert\vec F_{\mathrm{prop}}\rVert & total thrust-vector magnitude.\\
  ep1    & \epsilon_1 & thrust vector angle 1.\\
  ep2    & \epsilon_2 & thrust vector angle 2.\\
\end{tabularx}
\end{center}
If more than one propulsive force vector is supplied, \statevar{thrust} is the sum
of the individual magnitudes. That total has a direct physical interpretation only
when all thrust vectors are aligned. The angle outputs \statevar{ep1} and
\statevar{ep2} correspond to the last thrust vector evaluated; multiple
\problemblock{*prop} data blocks are evaluated in input order.
